\hypertarget{sec:hipotesis}{}%
\begin{center}
    {\Large\bfseries 5. FORMULACIÓN DE HIPÓTESIS}
\end{center}
\vspace{1em}

\noindent Dado que el alcance del estudio es descriptivo con un componente comparativo secundario (numeral 4), las hipótesis se formulan de dos tipos: \textit{descriptivas} (H1a y H1b), que anticipan la tendencia de una variable sin relacionarla con otra, y de \textit{diferencia de grupos} (H1c), que anticipa que la afectación difiere entre las aplicaciones. No se formulan hipótesis causales, pues el diseño no manipula ni controla las caídas del servicio. Cada hipótesis de trabajo (H1) se acompaña de su hipótesis nula (H0), que es la que se pone a prueba en el análisis de resultados: no se demuestra H1 de forma directa, sino que se evalúa si existe evidencia suficiente para rechazar H0. Los objetivos específicos a) y e), de caracterización y de recomendaciones, no requieren hipótesis propia.
\par\vspace{1.2em}

\noindent Las hipótesis no fijan un porcentaje, porque los antecedentes disponibles no permiten justificar un valor de corte; la proporción observada se estimará y se reportará en los resultados. Las hipótesis se refieren a los participantes de la muestra; cualquier generalización a una población más amplia deberá definirse en el diseño metodológico, junto con la población a la que se refiere y los supuestos que la sostienen. Se redactan de modo que un resultado sin diferencia, o en sentido contrario, también constituya un hallazgo válido.
\par\vspace{1.2em}

\noindent\textbf{H1a (hipótesis de trabajo --- nivel de uso):} En la muestra predominan los estudiantes de pregrado de la UPTC, Sede Tunja, que realizan la mayoría o casi todos los pagos de sus actividades cotidianas (alimentación, transporte, fotocopias y pagos académicos) con al menos una de las tres aplicaciones (Nequi, DaviPlata o Bancolombia).
\par\vspace{0.8em}

\noindent\textbf{H0a (hipótesis nula --- nivel de uso):} En la muestra no predominan los estudiantes que realizan la mayoría o casi todos los pagos de sus actividades cotidianas con al menos una de las tres aplicaciones.
\par\vspace{1.2em}

\noindent\textbf{H1b (hipótesis de trabajo --- consecuencias de la caída):} Entre los estudiantes de la muestra que vivieron en el semestre una caída o falla de alguna de las tres aplicaciones, confirmada como problema de la aplicación y no de su conexión o dispositivo, predominan quienes reportan al menos un tipo de afectación en alguna actividad cotidiana (aplazar o retrasar la actividad, cancelar una compra o un pago, no cumplir una actividad académica a tiempo o tener un costo adicional).
\par\vspace{0.8em}

\noindent\textbf{H0b (hipótesis nula --- consecuencias de la caída):} Entre los estudiantes de la muestra que vivieron una caída o falla confirmada, no predominan quienes reportan al menos un tipo de afectación en alguna actividad cotidiana.
\par\vspace{1.2em}

\noindent\textbf{H1c (hipótesis de trabajo --- comparación secundaria entre aplicaciones):} Entre los estudiantes de la muestra que usan cada aplicación, la afectación de las actividades financieras por su no disponibilidad difiere según la aplicación (Nequi, DaviPlata o Bancolombia).
\par\vspace{0.8em}

\noindent\textbf{H0c (hipótesis nula --- comparación secundaria entre aplicaciones):} Entre los estudiantes de la muestra que usan cada aplicación, la afectación de las actividades financieras por su no disponibilidad no difiere según la aplicación.
\par\vspace{1.2em}

\noindent La Tabla~\ref{tab:hipotesis-resumen} sintetiza la relación entre cada par de hipótesis, su tipo, el objetivo específico que responde y las variables que involucra.

\begin{table}[!ht]
\centering
\small
\renewcommand{\arraystretch}{1.3}
\begin{tabular}{|c|p{3.4cm}|p{3.6cm}|p{4cm}|}
\hline
\textbf{Hipótesis} & \textbf{Tipo} & \textbf{Objetivo} & \textbf{Variable(s) involucrada(s)} \\ \hline
H1a / H0a & Descriptiva & b) Nivel de uso & Uso \\ \hline
H1b / H0b & Descriptiva & c) Consecuencias de la no disponibilidad & Falla reportada y afectación \\ \hline
H1c / H0c & Diferencia de grupos (comparativa secundaria) & d) Cuantificar la afectación & Aplicación y afectación \\ \hline
\end{tabular}
\caption{Síntesis de las hipótesis de la investigación.}
\label{tab:hipotesis-resumen}
\end{table}
